% =====================================================================
%  CAPÍTULO 5 — ARQUITETURA E DESIGN DA SOLUÇÃO
% =====================================================================
\chapter{Arquitetura e Design da Solução}
\label{cap:arquitetura}

\section{Visão Global da Arquitetura}
\label{sec:visao_global}

O sistema segue uma arquitetura \textbf{cliente-servidor} de três camadas
com um \emph{pipeline} de processamento de quatro estágios no núcleo do
servidor. O cliente é uma \emph{Single-Page Application} (\textsc{spa})
React; o servidor é uma API \textsc{rest} assíncrona em FastAPI; e a
camada de dados é externalizada para serviços geridos do Supabase
(autenticação, base de dados relacional e armazenamento de objetos). A
\Cref{fig:arquitetura} (Anexo~\ref{ap:diagramas}) representa a organização
global e os fluxos principais.

Uma característica determinante da arquitetura é o modo como se gerem as
operações demoradas dos módulos M1--M4 (\textsc{ocr}, análise visual,
inferência do \textsc{llm} e montagem de vídeo), que podem demorar de
segundos a minutos. O sistema dispõe de uma \textbf{infraestrutura de
tarefas em segundo plano} --- com duas vias intermutáveis (ver
\Cref{sub:async}): uma \textbf{fila de tarefas} (Celery + Redis) quando
existe um \emph{worker} dedicado, ou um \textbf{executor \emph{in-process}}
dentro do próprio \emph{backend} quando não existe (o caso do plano
gratuito) ---, na qual o utilizador recebe de imediato um identificador de
tarefa e o \emph{frontend} acompanha o progresso por \emph{polling}. Esta
via é usada, em produção, para a \textbf{análise visual diferida das
fotografias} (M1). Em contrapartida, a \textbf{geração de narrativa e de
vídeo corre de forma síncrona}, mantendo a instância acordada durante o
processamento: uma escolha deliberada, motivada pela fragilidade dos
\emph{jobs} de segundo plano no plano gratuito (instâncias que adormecem ou
reiniciam, abandonando a tarefa --- \Cref{sub:async}). A
\Cref{fig:fullstack} (Anexo~\ref{ap:diagramas}) apresenta uma vista
detalhada, camada a camada, de toda a \emph{stack}.

\section{Evolução Arquitetural: de \emph{local-first} a multi-inquilino}
\label{sec:evolucao}

O projeto nasceu com uma filosofia \textbf{\emph{local-first}}: toda a
computação e persistência ocorriam na máquina do utilizador, com
autenticação \textsc{jwt} \textsc{hs256} própria, base de dados SQLite e
ficheiros guardados localmente. Este desenho cumpria o princípio de que
\emph{nenhum dado saísse da máquina}, e corresponde ao que foi inicialmente
planeado: SQLite como base de dados estruturada e o sistema de ficheiros
para os \emph{media} originais.

A necessidade de disponibilizar o sistema publicamente motivou uma
\textbf{migração para uma arquitetura de nuvem multi-inquilino}:

\begin{itemize}[leftmargin=*]
  \item A \textbf{autenticação} passou de um fluxo HS256 próprio para o
  \textbf{Supabase Auth}, com tokens assinados por chave ECC P-256 e
  validados contra o \emph{endpoint} JWKS público.
  \item O \textbf{modelo de domínio} deixou de ser de dono único
  (\emph{single-owner}) para \textbf{multi-inquilino}: todas as entidades
  referenciam um
  \texttt{user\_id} (\textsc{uuid}) e todas as consultas filtram por esse
  campo.
  \item O \textbf{armazenamento} migrou do disco local para o
  \textbf{Supabase Storage}, passando a referência de cada ficheiro a ser
  uma chave de objeto em vez de um caminho de disco.
\end{itemize}

É importante sublinhar que a migração para a nuvem \emph{não} eliminou o
modo local: o sistema corre de forma \textbf{idêntica nos dois ambientes}
--- tudo o que se faz \emph{online} é possível \emph{localmente}, e
vice-versa --- partilhando o mesmo código, a mesma base de dados e o mesmo
\emph{Storage} (Supabase). Esta paridade é um requisito explícito do
projeto e tem consequências concretas: por exemplo, os vídeos são
\emph{produzidos} em ambiente local (onde não há limite de memória) e
\emph{reproduzidos} no \emph{site} implantado, por escreverem no mesmo
\emph{Storage} (\Cref{sec:impl_m4}). Quanto à IA, o motor de texto por
omissão é o \textbf{Groq} (nuvem, gratuito e rápido), mas a inferência
\textbf{local} via \textbf{Ollama} permanece disponível como opção de
privacidade e de operação \emph{offline}. Esta evolução é um exemplo
concreto de compromisso maduro entre acessibilidade, privacidade e
escalabilidade.

\section{Fluxo de Dados e Execução}
\label{sec:fluxo}

\subsection{Encadeamento do \emph{pipeline}}

Os quatro módulos formam uma cadeia em que a saída de cada um alimenta o
seguinte. A \Cref{tab:fluxo} resume, para cada módulo, o que consome e o
que produz — tornando explícito porque é que a ordem importa (por exemplo,
o M3 só pode ancorar a narrativa em factos depois de o M1 os ter
extraído).

\begin{table}[H]
\centering
\caption{Fluxo de dados ao longo do \emph{pipeline} M1--M4.}
\label{tab:fluxo}
\small
\begin{tabularx}{\linewidth}{@{}p{1.5cm}YY@{}}
\toprule
\textbf{Módulo} & \textbf{Consome} & \textbf{Produz} \\
\midrule
M1 & Bytes de fotografias, documentos e GEDCOM & \texttt{MediaFile} (metadados
\textsc{exif} + descrição de IA + \textsc{ocr}); \texttt{Person} e relações. \\
M2 & \texttt{MediaFile} e \texttt{Person} & Datas resolvidas;
\texttt{TimelineEvent}; grafo de parentesco. \\
M3 & Factos do M1/M2 (via RAG) + intenção do utilizador & \texttt{Story}
(prosa + \texttt{scenes} ligadas a fotografias). \\
M4 & \texttt{Story} (cenas) + fotografias do \emph{Storage} &
\texttt{VideoOutput} (MP4 sincronizado, no \emph{Storage}). \\
\bottomrule
\end{tabularx}
\end{table}

\subsection{Sequência de uma tarefa em segundo plano}

Quando uma operação corre em segundo plano (em produção, a análise visual
diferida das fotografias do M1), o fluxo é o seguinte:

\begin{enumerate}[leftmargin=*]
  \item A API valida o pedido, cria um registo \texttt{TaskRecord} (estado
  \texttt{pending}) e devolve \emph{imediatamente} o seu identificador.
  \item A tarefa é despachada para o Celery ou para o executor
  \emph{in-process} (\Cref{sub:async}), que a marca \texttt{running}.
  \item O corpo da tarefa executa; ao terminar, o registo passa a
  \texttt{done} (ou \texttt{failed}) e é ligado ao artefacto produzido.
  \item O \emph{frontend} faz \emph{polling} do estado e atualiza a
  interface na transição para \texttt{done}.
\end{enumerate}

A geração de \textbf{narrativa} e de \textbf{vídeo} pode, em código, ser
pedida nesta modalidade (\texttt{mode=background}); contudo, na
configuração implantada corre \textbf{sincronamente} (\Cref{sub:async},
\Cref{sec:impl_m4}), mantendo a instância acordada durante o processamento
e devolvendo o resultado diretamente --- evitando que uma instância gratuita
que adormeça abandone a tarefa a meio.

\section{Modelo de Dados}
\label{sec:modelo_dados}

A camada de persistência relacional é modelada com SQLAlchemy e materializada
em PostgreSQL (Supabase). As entidades principais são:

\begin{itemize}[leftmargin=*]
  \item \texttt{MediaFile} — cada ficheiro carregado (fotografia,
  documento, vídeo). Guarda o nome original, a chave de objeto no Storage,
  os metadados \textsc{exif} (data, GPS, câmara), os campos de análise de
  IA (\texttt{ai\_description}, \texttt{ai\_setting}, \texttt{ai\_emotion},
  \texttt{ai\_tags}, \texttt{ai\_narrative\_hint}), o texto de \textsc{ocr},
  o estado de processamento e \texttt{person\_ids} (as pessoas da árvore
  etiquetadas na fotografia).
  \item \texttt{Person} — cada pessoa, importada do GEDCOM \emph{ou criada
  manualmente} (nome, sexo, datas, locais, identificador GEDCOM, etiqueta
  de família).
  \item \texttt{Relationship} — uma aresta de parentesco entre duas
  pessoas (\texttt{pai}/\texttt{mãe}/\texttt{cônjuge}). As relações deixaram
  de viver apenas num grafo JSON em disco (efémero) e passaram a ser
  persistidas na base de dados, sustentando o editor manual e a vista
  interativa da árvore.
  \item \texttt{TimelineEvent} — eventos da linha temporal (título,
  descrição, tipo, local, data resolvida).
  \item \texttt{Story} — narrativa gerada (título, texto, tipo de evento,
  \emph{template} usado, \emph{backend} de \textsc{llm}, número de factos
  usados, \emph{prompt}, idioma, projeto associado). Guarda ainda as
  \emph{preferências de produção} do vídeo --- voz (\texttt{voice}),
  legendas (\texttt{subtitles}, \texttt{subtitle\_size}) ---, a lista de
  fotografias explicitamente escolhidas (\texttt{media\_ids}), um indicador
  de \emph{favorita} (\texttt{favorite}) e o campo \texttt{scenes}
  (\textsc{json}): a forma segmentada em cenas da narrativa --- uma lista de
  \texttt{\{text, photo\_ids, caption\}} que permite ao M4 sincronizar cada
  fotografia com a sua narração (\Cref{sub:cenas}).
  \item \texttt{VideoOutput} — vídeo gerado (nome, chave de objeto,
  tamanho, fotografias usadas, estado, história e projeto associados), bem
  como a referência à \emph{faixa de legendas} \texttt{.vtt} e o tamanho de
  legenda escolhido, consumidos pelo leitor (\Cref{sub:narracao_legendas}).
  \item \texttt{Project} e \texttt{ProjectMedia} — projetos (coleções) e
  a relação muitos-para-muitos com as fotografias. Cada projeto é um
  espaço de trabalho \emph{isolado}: as suas fotografias, a sua linha
  temporal (derivada apenas dessas fotografias), a sua família/árvore
  (identificada pela etiqueta do projeto), as suas histórias e vídeos não
  se misturam com os de outros projetos nem com a Biblioteca global.
  \item \texttt{TaskRecord} — registo de tarefas assíncronas (estado,
  \emph{payload}, carimbos temporais, erro). Cada registo guarda também o
  \texttt{project\_id} a que pertence, de modo que as tarefas em segundo
  plano de um projeto não se misturam com as dos restantes (filtragem por
  projeto na página de Tarefas).
\end{itemize}

As relações principais entre entidades resumem-se na \Cref{tab:relacoes}.
Uma fronteira de tipos relevante: o \texttt{user\_id} é um \textsc{uuid}
(proveniente do Supabase Auth), ao passo que as chaves primárias das
entidades de domínio são inteiros sequenciais do PostgreSQL. A
\Cref{fig:er} (Anexo~\ref{ap:diagramas}) apresenta o diagrama
entidade-relação.

\begin{table}[H]
\centering
\caption{Relações principais entre entidades.}
\label{tab:relacoes}
\small
\begin{tabular}{@{}lcl@{}}
\toprule
\textbf{Entidade A} & \textbf{Cardinalidade} & \textbf{Entidade B} \\
\midrule
Project      & 1 --- N & ProjectMedia \\
ProjectMedia & N --- 1 & MediaFile \\
Story        & 1 --- N & VideoOutput \\
Project      & 1 --- N & Story \\
TaskRecord   & N --- 1 & Story / VideoOutput (resultado) \\
\bottomrule
\end{tabular}
\end{table}


\section{Design de Interface (UI/UX)}
\label{sec:uiux}

\subsection{Princípios e sistema de design}

O \emph{design} visual procura evocar o \textbf{calor do arquivo
familiar}. Assenta numa paleta \emph{amber}/\emph{stone} (tons quentes
que remetem para fotografias antigas e luz de candeeiro), com acentos
secundários sóbrios. A tipografia é dupla: \emph{Inter} para a interface
(legibilidade) e \emph{Fraunces} (serifada) para títulos e para a prosa
do leitor de histórias, conferindo um registo literário. As animações são
discretas (\emph{fade-in}, \emph{slide-up}, \emph{skeletons} de
carregamento), realizadas com animações \textsc{css} declaradas no Tailwind.

Um detalhe de implementação relevante para a experiência é o
\textbf{tema} (claro/escuro/sistema): um \emph{script} embutido no
cabeçalho lê a preferência guardada \emph{antes} do primeiro
\emph{paint}, evitando o desagradável ``flash'' de tema errado ao
carregar a página.

\subsection{Estrutura de navegação}

A aplicação organiza-se em rotas \textbf{públicas} (\emph{landing}/login,
registo, recuperação de palavra-passe, ``sobre'') e \textbf{protegidas},
acessíveis após autenticação através de uma barra lateral
(\emph{sidebar}) responsiva. As secções visíveis são: \emph{Início},
\emph{Biblioteca}, \emph{Família}, \emph{Linha temporal}, \emph{Projetos},
\emph{Gerar}, \emph{Histórias}, \emph{Vídeos} e \emph{Definições}. A página
de \emph{Tarefas} permanece acessível pela rota \texttt{/tasks}, mas foi
\textbf{retirada da barra lateral} por as gerações passarem a correr de
forma síncrona, tornando o acompanhamento por \emph{polling} dispensável no
uso corrente.

As capturas de ecrã dos ecrãs centrais da aplicação --- \emph{Dashboard},
Biblioteca, assistente \emph{Gerar}, leitor de histórias e galeria de
vídeos --- constam do Anexo~\ref{ap:capturas}.

\subsection{Internacionalização e acessibilidade}

Toda a interface está disponível em \textbf{português europeu} (omissão)
e \textbf{inglês}, com a preferência persistida e o atributo
\texttt{<html lang>} sincronizado. Ao nível da acessibilidade, existem
rótulos \textsc{aria} e navegação por teclado (por exemplo, no
\emph{lightbox} da biblioteca); uma auditoria formal de conformidade
\textsc{wcag} é apontada como trabalho futuro (\Cref{cap:conclusoes}).
